\section{Implementation}
\label{sec:implementation}

We implement LatchMoE as a Python package that registers through vLLM's general
plugin entry point. The package integrates with a hook-enabled vLLM-Ascend
runtime; the upstream request scheduler and OpenAI-compatible serving API
remain unchanged. The host runtime and the Ascend platform plugin are pinned as
a compatible stack because the staging boundary changes the internal fused-MoE
execution contract rather than the public serving interface.

\subsection{Graph integration}

The integration decomposes a managed MoE layer into a captured router, the
custom \texttt{vllm::moe\_offload\_stage} splitting operation, and a
capacity-selected expert path. The splitting operation returns control to the
host after routing so that LatchMoE can stage the active experts. A fitting
invocation continues into the captured expert MLP, which reads persistent slot
weights and the logical-to-slot buffer. An overflowing prefill instead invokes
the eager wave executor against transient stage banks. The runtime locks the
persistent buffers' data pointers at capture and validates them before replay.

Graph-compatible runs use PIECEWISE ACLGraph. The benchmark validator requires
explicit capture and replay records for graph-designated pieces and rejects
forced eager execution or unintended eager fallback of those pieces, graph
errors, address changes, stale H2D completion, and slot-version violations. The
planned eager overflow path is not classified as fallback. Disabling LatchMoE
restores the host runtime's null hooks. The
launcher also rejects simultaneous use of the native weight-offload path and
LatchMoE because two independent owners of expert storage would violate the
slot lifecycle.

\subsection{Loading and memory management}

CPU-first loading creates managed expert tensors in host memory during model
initialization. The host store keeps one contiguous $W_{13}$ tensor and one
contiguous $W_2$ tensor per layer and exposes per-expert views to the transfer
engine. It checks shape, dtype, stride, device, layer coverage, and expert
coverage before releasing the original device expert banks. Host pinning is
used when the runtime supports it and reports failures explicitly.

Automatic configuration derives the resident-layer selection and slot count
from the requested offload amount, physical HBM, a KV-cache reserve, and a cap
on the fraction of HBM assigned to slot storage. Explicit overrides remain
available for controlled sensitivity experiments. The memory ledger accounts
separately for the host store, persistent slots, transient prefill stage banks,
mapping tensors, optional full-layer staging storage, and any original weights
that remain allocated.

\subsection{Transfers and lifecycle enforcement}

The transfer engine maintains a dedicated H2D stream. It validates layout
compatibility before copying and batches consecutive host and slot slices when
their storage is contiguous. Each asynchronous load returns a readiness object
that contains the NPU event and the exact slot leases made consumable by that
event. The consumer installs the event dependency before the object marks the
leases ready. The runtime then publishes the active persistent mapping or
wave-local binding.

Compute protection uses an event recorded after the expert operation. Slots
remain in the computing state until the event completes and every recorded
lease still matches its current owner and version. Shutdown drains pending
transfer and compute work, closes profiling writers, releases managed storage,
and emits a release acknowledgement for the benchmark harness.

\subsection{Wave execution and profiling}

The prefill path supports a configurable number of stage banks and prefetch
depth. The qualified configuration uses two stage banks and depth one. Wave
plans carry the original pair offsets needed for one native recombination, and
strict validation rejects duplicate or missing coverage. A controlled
evaluation configuration may replace bounded waves with explicit full-layer
staging; it is not an error-recovery path. NPU, ACL, memory, address, semantic,
and lifecycle failures are not swallowed by either execution mode.

Profiling records JSONL events for layer registration, slot ownership,
transfers, mapping publication, wave staging and compute, graph replay issue,
memory accounting, and service release. Benchmark manifests bind each run to
the runtime revisions, model and workload digests, device, command line,
environment, graph mode, and raw client/server artifacts. These records support
the fail-closed checks used in Section~\ref{sec:evaluation}.
